\pdfoutput=1
\documentclass[11pt]{article}

\usepackage{ACL2023}
\usepackage{times}
\usepackage{latexsym}
\usepackage[T1]{fontenc}

\usepackage[utf8]{inputenc}
\usepackage{microtype}
\usepackage{listings}
\usepackage{inconsolata}
\usepackage{fancyhdr}
\usepackage{cleveref}
\usepackage{natbib}

\pagestyle{fancy}
\fancyhf{}
\fancyhead[LE,RO]{EPFL – Spring Semester 2026}
\fancyhead[RE,LO]{Modern Natural Language Processing - Standard Project Final Report}
\rfoot{Page \thepage \hspace{1pt}}
\renewcommand{\headrulewidth}{1pt}

%%%%%%%%%%%%%%%%%%% TITLE AND AUTHORS %%%%%%%%%%%%%%%%%%%

\title{Post-training Qwen3-1.7B for 5 reasoning tasks}

\author{\normalfont 
Léonard Amsler | 346328 | \texttt{leonard.amsler@epfl.ch} \\
Noé Boulud | 326805 | \texttt{noe.boulud@epfl.ch} \\
Nathan Gromb | 341157 | \texttt{nathan.gromb@epfl.ch} \\
Nicolas Teissier | 345506 | \texttt{nicolas.teissier@epfl.ch} \\
4neurons \\
}

%%%%%%%%%%%%%%%%%%% COMMANDS %%%%%%%%%%%%%%%%%%%

\usepackage{tcolorbox}
\tcbuselibrary{breakable, skins}

\newtcolorbox{promptbox}{
  breakable,
  enhanced,
  colback=gray!10,
  colframe=gray!50,
  arc=3pt,
  boxrule=0.5pt,
  left=6pt, right=6pt, top=4pt, bottom=4pt,
  fontupper=\small\ttfamily
}

%%%%%%%%%%%%%%%%%%% REPORT %%%%%%%%%%%%%%%%%%%

\begin{document}
\maketitle
\thispagestyle{fancy}


\begin{abstract}
%A concise summary (under 300 words). State the gaps or motivation behind the work, your approach to producing the individual and group models, and the main empirical findings. Keep implementation details out of the abstract — those belong in the body.
\end{abstract}





\section{Introduction}
%Motivate the problem. Sketch how current methods fall short (with references to prior work), give a high-level overview of your individual and group model strategies, and preview the headline results. Example citation: \cite{Gusfield:97}.






\section{Method}
\label{sec:method}
%Describe your training approach in enough detail that a reader could re-implement it. Cover all five models — the four individual models (one per benchmark) and the group model. For each, motivate the choice of post-training method (SFT, RLHF, RLVR, or a combination), formalize the training objective where it helps, and explain non-trivial implementation choices. Clarify how the group model is obtained from the individual models, distillation, model merging, or just how you balanced specialization against generalization. Use equations and figures wherever they aid clarity. 
%You can create separate sections like the following:

\subsection{Data}

\subsubsection{Maths}

    The abundance of public mathematical datasets allows a light pre-processing. Some of the most popular dataset already contain distilled reasoning traces and answer verifications. An advantage of this dataset diversity is the ability to mix them to create large, diverse and general pools, with filtering of duplicates. However, as traces are generated with heavy state-of-the-art models, these are filtered until the provided limit of 16k tokens to avoid future generation of these overlong answers that will be left unfinished without a box, worsening performances. For multi-generations cases, the smallest one is chosen to encourage conciseness. Beside training sources, untouched benchmarks are used to evaluate out-of-distribution performances of the produced models. 

    
\subsubsection{General Knowledge}

    General knowledge is framed as multiple-choice questions with $2$ to $20$ options; the output contract is a final answer in \texttt{\textbackslash boxed\{LETTER\}}, with thinking baked into the chat template at merge time. We build the training pool from eight public MC sources spanning five macro-categories (STEM, humanities, social sciences, history/geography, commonsense): MMLU, MMLU-Pro, a world-knowledge MMLU subject mix, TriviaQA (reformatted to MC), BoolQ, CommonsenseQA, ECQA and SocialIQA. A key lesson from early iterations is that placeholder or synthetic reasoning traces teach the model \emph{not} to reason, so every retained example carries a genuine chain-of-thought (CoT) ending in \texttt{\textbackslash boxed\{\}}. Augmentation is controlled (per-macro-category typed distractor pools, option expansion up to $20$, capped at $\approx$20--30\%) so that the option-count distribution covers the $2$--$20$ regime without injecting absurd distractors. We compare two ways of producing the CoTs (analysed in the post-training section): \textbf{off-policy distillation} from Qwen3-14B-AWQ (short traces for fact-lookup sources, long step-by-step traces for STEM), and \textbf{on-policy self-distillation} from Qwen3-1.7B itself, keeping only its correct draws and filtering degenerate loops.


\subsubsection{Multilinguality}
Due to the possibility of having to choose between 2 and 20 MCQ choices, the very first step the multilingual data processing is to use a larger model to generate additional MCQ choices, called "distractors". The model is instructed to generate plausible distractors, that are incorrect and do not overlap with existing ones (see the prompt in \cref{appendix:multilingual_distractor_prompt}). 

\begin{figure}[h]
    \centering
    \includegraphics[width=\linewidth]{assets/questions_distribution.png}
    \caption{Distribution of the number of MCQ choices per question.}
    \label{fig:distractor_dist}
\end{figure}

We then sample between 2 and 20 choices at random to keep as our final ones, always including the right answer and giving priority to original MCQ choices to keep the quality as high as possible. The sampling process is designed to obtain many questions with 4 choices, and some examples of higher and lower choice counts following a Gamma-Uniform mixture, as illustrated in \cref{fig:distractor_dist}
    
The second step consists in using a larger model, which is better at multilingual tasks, and having it answer questions from the dataset as we would with our own Qwen-1.7B. We do not give the large model any hints, and filter based on the correctness of its final answers. The resulting distilled Chain-of-Thought traces are used as training data to improve the reasoning of our model.


\subsubsection{Safety}
\begin{enumerate}
    \item generate distilled CoT for MCQs questions on safety using Qwen3-32B as the "teacher" model, 8 tries, keep only right ones ($\approx$ 9'500 samples out of $\approx$ 10'200 samples selected for training)
    \item data divided into subcategories (ex: Offensiveness, Illegal Activities, Physical Health, ...), not given to the teacher model, only for results analysis purpose. Was already divided within the dataset (cf Dataset section maybe ?).
\end{enumerate}


\subsubsection{Group Model}
Merge datasets (details for maths)


\subsection{Post-training}

\subsubsection{Maths}

    To increase performances of the model on mathematical thinking, several approachs are explored:
    
    \begin{itemize}
        \item Prompt engineering: Exploration of the impact of three prompting strategies: no prompt, simple prompt for box formatting and enhanced math prompt.
        \item SFT: Training of a LoRA on a mixture of datasets. Several training configuration are tested including duration, dataset portions and maximum length of generation samples. 
        \item GRPO: Update the model with the GRPO algorithm from generations of the training datasets and others.
    \end{itemize}

    
\subsubsection{General Knowledge}

    Every variant is a LoRA fine-tune of Qwen3-1.7B, optionally followed by preference optimisation; the merge step bakes \texttt{enable\_thinking=ON} into the chat template. Rather than a single recipe, we ran a controlled sequence that isolates one factor at a time. V6 is our strongest purely off-policy model, supervised on 14B CoTs (long step-by-step traces on STEM, short ones elsewhere). We then compare two on-policy SFT variants on identical self-distilled data, V8 ($r{=}8$) and V9 ($r{=}64,\ \alpha{=}128$), to isolate the effect of LoRA capacity on the \texttt{\textbackslash boxed\{\}} format, and add V9b, which at assembly keeps the shortest \emph{clean} correct trace per question (falling back to the V6 CoT when all draws loop) to counter the 1.7B's tendency to repeat. V10 returns to off-policy with contrastive 14B CoTs that justify the gold answer and refute the most tempting wrong line, kept option-agnostic so they stay reusable across augmented variants. Finally, V11 and V11b apply on-policy DPO on top of V9b: we sample $8$ completions per prompt, form (correct, incorrect) preference pairs sharing the same format, and optimise the DPO objective with reference policy $\pi_{\text{ref}}=$ V9b, timidly for V11 ($\beta{=}0.1$, lr $5\mathrm{e}{-}6$, $1$ epoch) and with a stronger signal for V11b ($\beta{=}0.05$, lr $1\mathrm{e}{-}5$, $3$ epochs, \texttt{load\_best\_model\_at\_end}).


\subsubsection{Multilinguality}
SFT, SFT with grouped batches, GRPO


\subsubsection{Safety}

\begin{enumerate}
    \item SFT (LoRA) on the CoT augmented data, multiple epochs
    \item for DPO focus only on weak categories (all samples of categories with weak performances measured, cf evaluation section), make the (post SFT) model generate 8 answers, take first correct answer as chosen, first wrong answer as rejected, and train DPO on this
\end{enumerate}


\subsubsection{Group Model}
%\begin{enumerate}
    %\item LoRA merging techniques (CAT (static/learnable), TIES, DARE-TIES TODO:cite papers), hyperparameters optimized to achieve best results on validation split composed of mix of validation split of each domain specific dataset (cf section datasets maybe ?)
    %\item SFT with grouped batches (one training dataset present per batch to avoid instability in training TODO: find paper supporting this)
%\end{enumerate}

We explore two distinct categories of approaches for merging the capacities of our four specific models into a group model : 
\begin{enumerate}
    \item LoRA weight merging techniques for skill composition taken from \citeauthor{prabhakar2024lorasoupsmergingloras}
    \item Supervised fine-tuning using LoRA 
\end{enumerate}

The clear advantage of the former is its computational simplicity: while some parameters may need to be optimised, for example to weight the layers of the different LoRA adapters that are being merged, the amount of compute is orders of magnitude lower \{SOURCE\} than that of SFT or even PEFT. This allows us to reuse the results of our different task-specific experiments. Among the existing methods, we run \{NICO? CAT (static/learnable), TIES, DARE-TIES\} and compare their performances.

In parallel, we also run LoRA fine-tuning on a mixture of the four task-specific datasets. In order to increase performance, we take inspiration from \citeauthor{rao-etal-2024-commonit} and use task-homogeneous batching.

\section{Experimental Setup}


\subsection{Datasets}
%For each model, describe the training data: sources, sizes, preprocessing steps, and any filtering or augmentation. Note licensing, legal, or ethical considerations. Check if a summary table aids clarity for the reader.

\subsubsection{Math}

    Below is the list of all preexisting datasets used for training on mathematical thinking:

    \begin{enumerate}
        \item OpenMathInstruct-2 (OMI) [CITE]: 14M rows from solution and problem-solution augmentations of the MATH and GSM8K with reasonning traces of the Llama3.1-405B-Instruct model. Used a portion the train1M split.
        \item OpenR1-Math-220k (OR1M) [CITE]: 
        \item NuminaMath-1.5 [CITE]:
        \item MATH500 [CITE]: Benchmark, leakage verification of with OMI, OR1M and NuminaMath has been made with word and embeddings overlap. 
    \end{enumerate}
    
    For SFT, the data mix includes OMI and OR1M. For GRPO, NuminaMath is added as no traces are needed.

\subsubsection{Safety}

\begin{enumerate}
    \item SafetyBench: "a compre-
    hensive benchmark for evaluating the safety
    of LLMs, which comprises 11,435 diverse
    multiple choice questions spanning across 7
    distinct categories of safety concerns." TODO: cite paper 
    \item divided into train and validation split (test split is the CI pipeline on hugging face)
    \item motivation for safetybench -> proven that "safety understanding abilities in SafetyBench
    are correlated with safety generation abilities", right format (MCQ), spanning different safety categories (useful for analysis), amount of data sufficient for training and validation)
    \item augmented through CoT generation (cf method/data section)
    \item table of data stats before CoT vs After CoT gen (by categories, number of tokens ?)
    
    
\end{enumerate}

\subsubsection{General Knowledge}

    The SFT mixture contains $\approx$30k examples drawn from the eight sources above and rebalanced across the five macro-categories (quotas $\approx$ STEM 25\%, humanities/social/history 20\% each, commonsense 15\%). Originals are deduplicated and a development blocklist (the three GK validation snapshots) is excluded before sampling. Roughly $20$--$30\%$ of rows are augmented variants (option shuffling / expansion with typed distractors). For the off-policy line, $\approx$139k CoTs were distilled from Qwen3-14B-AWQ (acceptance $\approx$96\%); for the on-policy line, the 1.7B baseline generated $4$ completions per prompt (thinking on, $T{=}0.6$), of which only correct draws were kept ($\approx$90\% of questions). All sources are public research MC datasets used under their respective licenses; questions overlapping any validation set are removed to prevent leakage.

\subsection{Evaluation}
%Specify the metrics you report, the validation sets (their construction, coverage, etc), and the baselines you compare against. Justify why each baseline is the right one to beat.

\subsubsection{Math}

    We evaluate model performances on pass@1 but more importantly pass@8 metrics on the MATH500 benchmark [CITE].

    Validation metrics were utilized during training to avoid overfitting but not as a performance metrics as it comes from the exact same distribution and could threaten generalization. 

    As baselines, we use the Qwen3-1.7B [CITE] model with no prompt, a simple prompt for the boxed formatting and an enhanced math one. As the models are trained and generate with a prompt, the latter make more sense to compare the model itself. 


\subsubsection{Safety}

\begin{enumerate}
    \item models' performance measured through pass@1
    \item also measured boxed format compliance (percentage of answers containing a boxed{} block at the end)
    \item evaluated on validation set (10\% of SafetyBench)
    
\end{enumerate}

\subsubsection{General Knowledge}

    We report \texttt{pass@1} and \texttt{pass@8}, scoring every completion with the official knowledge extractor (it reads the last \texttt{\textbackslash boxed\{\}}; without one it falls back to a single isolated choice letter, otherwise the answer is counted wrong). This makes the \texttt{\textbackslash boxed\{\}} format a first-class metric, since a correct but unformatted answer is lost. Beyond the global score we break results down by \emph{source} (BoolQ, MMLU, CommonsenseQA, MMLU-Pro), \emph{macro-category}, and \emph{option count} ($2$/$4$/$5$/$6$--$10$), which act as our different benchmarks. We evaluate on an out-of-training \texttt{dev\_full} ($n{=}1000$, $\pm0.031$), kept blocklisted from training; an earlier \texttt{dev\_small} ($n{=}220$) is too small to separate close models ($\pm0.065$) and is reported only for continuity. The baseline to beat is the raw Qwen3-1.7B (thinking on). Because a single $n{=}1$ draw is noisy ($\pm\sim3$ pts), we only rank close models at $n{=}8$.

\subsection{Implementation Details}
%Report hyperparameters, optimizer settings, training duration, hardware used, and anything else needed for reproducibility. Tables might help here again for clarity.

\subsubsection{Math}

    As generations are sensible to decoding parameters (temperature, top-p, top-k), we perform a decoding parameters search. Over 20 combinations are tested to measure their impact on performance metrics. Specially, it is interesting to track the tradeoff between pass@1 and pass@8 performances if more/less diverse tokens are sampled.

\subsubsection{General Knowledge}

    SFT uses LoRA on Qwen3-1.7B with $r{=}64,\ \alpha{=}128$, $1$ epoch, lr $2\mathrm{e}{-}4$, effective batch $16$ (bf16, \texttt{max\_seq\_length}$=4096$); the format-ablation V8 instead uses $r{=}8,\ \alpha{=}16$, lr $5\mathrm{e}{-}5$. DPO (V11/V11b) reuses the $r{=}16$ adapter on top of the merged V9b policy. After training, \texttt{merge\_lora} fuses the adapter, bakes \texttt{enable\_thinking=ON} and writes the decoding defaults ($T{=}0.7$, top-$p\ 0.9$, top-$k\ 20$) used at inference. Distillation and training run on a single A100; the full reproduction pipeline (data $\rightarrow$ SFT/DPO $\rightarrow$ merge $\rightarrow$ eval) is scripted per version in \texttt{pipelines/run\_v*.sh}, and all checkpoints can be re-evaluated on every benchmark with \texttt{pipelines/eval\_all.sh}.

\section{Results}

\subsection{Main Results}
%Present quantitative results with tables and plots. Compare each individual model against baselines on its dedicated benchmark, and report the group model's performance across all four benchmarks side-by-side.

\begin{itemize}
    \item group model SFT vs LoRA merging
    \item group model on each benchmark ? needs to run though
    \item loss curves over epochs for all tasks
    \item 
\end{itemize}

\subsection{Maths}

\begin{itemize}
    \item Performances per categories over all versions
    \item Performances per level over all versions
    \item Performances on benchmark over all versions
\end{itemize}

\subsection{General Knowledge}

\Cref{tab:gk_sweep} sweeps the kept versions on \texttt{dev\_full} per source, and \cref{tab:gk_n8} gives the reliable $n{=}8$ verdict for the two finalists against the baseline. Three controlled comparisons explain the ranking. First, \emph{format is driven by LoRA strength}: on identical data, moving from $r{=}8$ (V8) to $r{=}64$ (V9) lifts \texttt{\textbackslash boxed\{\}} coverage from $51\%$ to $86\%$ and pass@1 from $0.39$ to $0.60$; the baseline's near-zero BoolQ is the same effect, an unboxed yes/no being scored wrong. Second, \emph{on-policy beats off-policy}: the cleaned self-distilled V9b beats both 14B-taught models (V6, V10) on every source, the student learning better from its own loop-free reasoning than from a teacher it cannot reproduce. Third, \emph{DPO needs enough signal}: timid V11 is a no-op (chosen and rejected differ only by the final letter), whereas the stronger V11b gains $+2.0$ pass@1 at $n{=}8$ on $4/4$ sources for a negligible $-0.9$ pass@8 cost, sharpening the distribution toward consistent answers. Finally, \texttt{dev\_small} ($n{=}220$) returned $0.60$ for V9, V9b and V10 alike, a false plateau that only \texttt{dev\_full} and $n{=}8$ scoring resolve. \textbf{V11b is our final model, with V9b as fallback.}

\begin{table}[h]
\centering
\small
\begin{tabular}{lccccc}
\hline
Model & Global & BoolQ & MMLU & M-Pro & CSQA \\
\hline
V6   & 0.541 & 0.547 & 0.532 & 0.520 & 0.600 \\
V9b  & 0.609 & 0.567 & 0.647 & 0.537 & \textbf{0.693} \\
V10  & 0.554 & 0.513 & 0.585 & 0.517 & 0.587 \\
V11  & 0.610 & 0.593 & 0.655 & 0.530 & 0.667 \\
V11b & \textbf{0.640} & \textbf{0.640} & \textbf{0.678} & \textbf{0.567} & 0.687 \\
\hline
\end{tabular}
\caption{\texttt{dev\_full} pass@1 by source ($n{=}1$ draw; trend only, $\pm\sim3$ pts). Refreshed at $n{=}8$ by \texttt{pipelines/eval\_all.sh}.}
\label{tab:gk_sweep}
\end{table}

\begin{table}[h]
\centering
\small
\begin{tabular}{lcccccc}
\hline
Model & p@1 & p@8 & BoolQ & CSQA & MMLU & M-Pro \\
\hline
Baseline & 0.458 & 0.603 & 0.016 & 0.613 & 0.598 & 0.415 \\
V9b      & 0.580 & \textbf{0.892} & 0.574 & 0.644 & 0.603 & 0.521 \\
V11b     & \textbf{0.600} & 0.883 & 0.584 & 0.663 & 0.619 & 0.550 \\
\hline
\end{tabular}
\caption{\texttt{dev\_full} reliable $n{=}8$ evaluation. CI pass@1 and the $n{=}8$ rows for V6/V9/V10/V11 are added once the reruns complete.}
\label{tab:gk_n8}
\end{table}

\subsection{Group Model}

\begin{figure}
    \centering
    \includegraphics[width=1\linewidth]{assets/group_model_results.png}
    \caption{faire ça en mieux}
    \label{fig:group_model_results}
\end{figure}

\subsection{Analysis}
%Go beyond reporting numbers. Discuss what worked and what did not like ablations of method components, comparisons across post-training methods, training instabilities, and unexpected failure modes. Null results are informative: when something did not work, say so, hypothesize why, and what you tried to tackle it. Close with the main takeaways the reader should remember.

\section{Ethical Considerations}
%Discuss the broader impact of your work: who benefits, who could be harmed, and whether any harms fall disproportionately on marginalized groups. Suggest mitigations for the harms you identify (like adaptation to lower-resource languages), and argue why your work is, on balance, a positive contribution despite the possible harms.

Benefits: searchers for wide overview about post training techniques, both on a domain-specific and across domains tasks. 

Harmed: ?

Mitigations for each harms:

Impact of the work despite harms:

\section{Conclusion}
%Summarize your main contributions and findings, state the primary limitations of your work honestly, and suggest concrete directions for future investigation.

Main contribution: improve model capabilities on specific domains of NLP as well as grouping the capabilities into a single model without deteriorating too much performances. 

Limitations: Compute and Time constrains

Directions: Try more hyperparameters search (ex: ablation study on LoRA adapters hyperparameters), 

% \section*{Acknowledgements}
% This document has been adapted for CS-552 by Beatriz Borges from the ACL 2023 template. This template was created by Jordan Boyd-Graber, Naoaki Okazaki, and Anna Rogers, based on the style files used for earlier ACL, EMNLP and NAACL proceedings, including those for
% EACL 2023 by Isabelle Augenstein and Andreas Vlachos,
% EMNLP 2022 by Yue Zhang, Ryan Cotterell and Lea Frermann,
% ACL 2020 by Steven Bethard, Ryan Cotterell and Rui Yan,
% ACL 2019 by Douwe Kiela and Ivan Vuli\'{c},
% NAACL 2019 by Stephanie Lukin and Alla Roskovskaya, 
% ACL 2018 by Shay Cohen, Kevin Gimpel, and Wei Lu, 
% NAACL 2018 by Margaret Mitchell and Stephanie Lukin,
% Bib\TeX{} suggestions for (NA)ACL 2017/2018 from Jason Eisner,
% ACL 2017 by Dan Gildea and Min-Yen Kan, NAACL 2017 by Margaret Mitchell, 
% ACL 2012 by Maggie Li and Michael White, 
% ACL 2010 by Jing-Shin Chang and Philipp Koehn, 
% ACL 2008 by Johanna D. Moore, Simone Teufel, James Allan, and Sadaoki Furui, 
% ACL 2005 by Hwee Tou Ng and Kemal Oflazer, 
% ACL 2002 by Eugene Charniak and Dekang Lin, 
% and earlier ACL and EACL formats written by several people, including
% John Chen, Henry S. Thompson and Donald Walker.
% Additional elements were taken from the formatting instructions of the \emph{International Joint Conference on Artificial Intelligence} and the \emph{Conference on Computer Vision and Pattern Recognition}.

%%%%%%%%%%%%%%%%%%% BIBLIOGRAPHY %%%%%%%%%%%%%%%%%%%

% Entries for the entire Anthology, followed by custom entries
\bibliography{custom}
\bibliographystyle{acl_natbib}

%%%%%%%%%%%%%%%%%%% APPENDIX %%%%%%%%%%%%%%%%%%%
\newpage
\appendix

\section{Appendix}
\label{sec:appendix}

%This does not count towards your 4-page limit. You can put additional results, tables, and figures here; for example, extended hyperparameter sweeps, failed experiments, qualitative generation examples, or prompt templates. Your main body should still be self-sufficient, i.e., anything important to your argument should go there, and you should assume that you will be graded based on the content of the main part of your paper only.
\subsection{Prompt templates}

\subsubsection{Simple format prompt}\label{appendix:math_simple_prompt}
\begin{promptbox}
Solve the problem. Put only the final answer at the very end in LaTeX boxed format: $\textbackslash \textbackslash boxed\{answer\}$. Do not write anything after the box.
Examples: $\textbackslash \textbackslash boxed\{42\}$ if the answer is 42, $\textbackslash \textbackslash boxed\{A\}$ if the answer is A, $\textbackslash \textbackslash boxed\{p - q\}$ if the answer is p-q.
\end{promptbox}

\subsubsection{Mathematical reasoning prompt}\label{appendix:math_prompt}
\begin{promptbox}
You are a rigorous competition-math solver. Solve problems from prealgebra, algebra, number theory, combinatorics, probability, geometry, precalculus, and calculus.

Use only information stated in the problem. Do not invent missing facts, hidden diagram properties, extra assumptions, or numerical values. 

Work carefully and verify each important step. Prefer exact symbolic reasoning over decimal approximations. Keep track of definitions, constraints, and edge cases.

For geometry, justify angle, length, arc, parallel, cyclic, and similarity claims before using them.
For counting and probability, avoid double-counting and check whether cases are disjoint and exhaustive.
For algebra, check domains and discard extraneous solutions.

If there are multiple plausible approaches, choose the most reliable one and complete it. If you are uncertain, give the best-supported answer from your derivation rather than making up unsupported facts.

Your response may include concise reasoning, but the final answer must appear at the very end in LaTeX boxed format: $\textbackslash \textbackslash boxed\{answer\}$. Do not write anything after the final $\textbackslash \textbackslash boxed\{answer\}$.
Examples: If the answer is 42, write $\textbackslash \textbackslash boxed\{42\}$, if the answer is p-q, write $\textbackslash \textbackslash boxed\{p-  q\}$.
For polynomials, the answer should first display the higher order and be reduced as much as possible. Example $\textbackslash \textbackslash boxed\{x^3 - 300x^2 + 30000x - 999900\}$ instead of $\textbackslash \textbackslash boxed\{- 999900 + x^3 - 300x^2 + 30000x\}$ (incorrect order) or $\textbackslash \textbackslash boxed\{x^3 - 300x^2 + 30000x - 1000000 + 100\}$ (not reduced)
\end{promptbox}

\subsubsection{Multilingual distractor generation}\label{appendix:multilingual_distractor_prompt}

\begin{promptbox}
\textbf{Task}

You are given a multiple-choice question with one correct answer and several wrong answers. 
Your task is to generate additional wrong answers that are plausible but incorrect. 
The generated additional answers should not be obviously incorrect or nonsensical, and should be similar in style and content to the original choices.
In any case, the generated wrong choices must not include the correct answer or repeat any of the original wrong choices.

\textbf{Language}

The question and choices can be in any of the following languages: Italian, Spanish, Chinese, Russian, Hindi.
Make sure to generate wrong choices that are in the same language as the question and original choices.

\textbf{Example}

Question: What is the capital of France?
Choices: Berlin, Madrid, Paris, Rome
Additional wrong choices, 3 examples:
['Lisbon', 'Vienna', 'Brussels']
\end{promptbox}

\subsubsection{Multilingual prompt}\label{appendix:multilingual_cot_prompt}

\begin{promptbox}
\textbf{Task Instructions}

You are a helpful multilingual reasoning assistant for multiple-choice questions.
For multiple-choice questions, reason step by step, then provide your final answer as a single letter inside \textbackslash boxed\{\}, for example: \textbackslash boxed\{A\}.
Do not include anything after the \textbackslash boxed\{\}.

\textbf{Language Instructions}

Questions will be in one of the following languages: Italian, Spanish, Russian, Hindi, Chinese. 
You should reason in English, but follow the instructions below:
Think carefully, do not translate the question while solving. 
Preserve the question in the target language so that you keep all details without adding noise. 
After you finish thinking, state your answer in fluent and coherent target language.
Preserve named entities, quoted spans, and key terms in the target language while generating the rest of the reasoning in English.
\end{promptbox}

\end{document}
